\section{模型假设}\label{sec:assume}

\begin{enumerate}[label=(\arabic*),leftmargin=2.2em,itemsep=2pt]
\item 药材为均质、各向同性的长圆柱，初始温度和含水率均匀。长度是半径的 12.5 倍，端面只占总表面积的约 7\%，忽略端面的传热传质，温度和含水率只随半径和时间变化。
\item 药材内部的水分迁移用有效扩散描述\cite{whitaker}，扩散系数取题给经验公式；水分在表面离开药材，内部不计相变。
\item 表面换热服从牛顿冷却定律，表面失水服从对流传质定律，系数 $h=\SI{25}{W/(m^2.K)}$、$h_m=\SI{8e-7}{m/s}$ 取自附录 2，各问沿用。题设的热边界不含蒸发吸热，第~\ref{sec:evap}~节放宽这一假设。
\item 烘房空气充分混合，温度和含水浓度只随时间变化：前 $\SI{4}{h}$ 按附件 1 线性插值；此后烘房处于恒温段，取最后 $\SI{1}{h}$ 数据的平均值。
\item 附件 1 中烘房的含水浓度理解为空气含湿量（kg 水/kg 干空气），据此折算的相对湿度在 60\%--82\% 之间。这一理解只用于第~\ref{sec:evap}~节的湿球温度计算，不影响四问的结果。
\item 问题二、三中药材尺寸不变；问题四中药材按附件 2 的半径均匀收缩，干物质随药材骨架一起移动，72 h 之后半径保持末值。
\end{enumerate}
